% BHU PYQ -- Auto-generated & error-corrected
\documentclass[11pt,a4paper]{article}

\usepackage[a4paper,top=2.5cm,bottom=2.5cm,left=2.8cm,right=2.8cm,headheight=14pt]{geometry}
\usepackage{amsmath,amssymb}
\usepackage[T1]{fontenc}
\usepackage[utf8]{inputenc}
\usepackage{lmodern}
\usepackage{microtype}
\usepackage{fancyhdr}
\usepackage{enumitem}
\usepackage{hyperref}
\usepackage{lastpage}
\usepackage{parskip}

\hypersetup{colorlinks=true,urlcolor=black,linkcolor=black,
  pdftitle={STB-201: Descriptive Statistics and Distribution Theory -- BHU PYQ},pdfauthor={Banaras Hindu University}}

\pagestyle{fancy}\fancyhf{}
\renewcommand{\headrulewidth}{0.4pt}\renewcommand{\footrulewidth}{0.4pt}
\fancyhead[L]{\small   \textit{Banaras Hindu University}}
\fancyhead[R]{\small   \textit{STB-201: Descriptive Statistics and Distribu}}
\fancyfoot[C]{\small Page  \thepage\ of \pageref{LastPage}}


\newlist{parts}{enumerate}{2}
\setlist[parts,1]{label=(\alph*),leftmargin=*,itemsep=5pt,topsep=3pt}
\setlist[parts,2]{label=(\roman*),leftmargin=*,itemsep=3pt,topsep=2pt}

\newcommand{\pts}[1]{\hfill   \textit{[\#1]}}


\begin{document}


\begin{center}
  {\large\bfseries BANARAS HINDU UNIVERSITY}\\[2pt]
  {\normalsize B.Sc. (Hons.) Semester II Examination, 2023-24}\\[6pt]
  \rule{0.8\linewidth}{0.5pt}\\[6pt]
  {\large\bfseries Paper No. STB-201: Descriptive Statistics and Distribution Theory}\\[4pt]
  {\normalsize Time: 3 Hours\hfill Maximum Marks: 70}
\end{center}
\rule{\linewidth}{0.4pt}
\smallskip



\noindent   \textit{Instructions: Answer   \textbf{five} questions including Question~1 which is compulsory. Symbols have their usual meanings.}
\bigskip

\medskip


\noindent   \textbf{Question 1.}

\begin{parts}
  \item Define the coefficient of correlation. Show that it is unaffected by a change of origin but may be affected by a change of scale. \pts{7}
  \item Derive the formula for the rank correlation coefficient. If $(x_i, y_i)$ denotes the rank of the $i$-th pair ($i = 1, 2, \dots, n$) with $x_i - y_i = d_i$ and $x_i + y_i = n + 1$, then show that the rank correlation between $X$ and $Y$ is $-1$. \pts{7}
\end{parts}

\medskip\rule{\linewidth}{0.2pt}

\medskip


\noindent   \textbf{Question 2.}

\begin{parts}
  \item Using the principle of least squares, obtain expressions for the intercept and slope of the regression line $Y = a + bX$. The regression equation between income ($X$) and expenditure ($Y$) is $Y = 500 + 50X$. If the ratio of the standard deviation of expenditure to the standard deviation of income is $60$, find the coefficient of correlation between income and expenditure. \pts{7}
  \item Derive the expression to measure the acute angle between two regression lines. For a certain dataset, $Y - 1.2X = 0$ and $X - 0.6Y = 0$ are the regression lines. If $ \theta$ is the acute angle, find the value of $ an  \theta$. \pts{7}
\end{parts}

\medskip\rule{\linewidth}{0.2pt}

\medskip


\noindent   \textbf{Question 3.} Explain the concept of multiple and partial correlations. In the case of three variables $X_1, X_2$, and $X_3$, show that:
$$1 - R_{1.23}^2 = \frac{|R|}{|R_{11}|}$$
where $|R|$ is the determinant of the correlation matrix of variables $X_1, X_2$, and $X_3$, and $|R_{11}|$ is the cofactor of $r_{11}$ in $|R|$. Justify that the range of $R_{1.23}$ is between zero and one. \pts{14}

\medskip\rule{\linewidth}{0.2pt}

\medskip


\noindent   \textbf{Question 4.}

\begin{parts}
  \item Three fair coins are tossed. Let $X$ denote the number of heads on the first two coins and $Y$ denote the number of tails on the last two coins. Find the joint distribution of $X$ and $Y$. What will be the value of $ \text{Cov}(X, Y)$? \pts{7}
  \item Define the moment generating function (MGF) of a random variable and write any three of its properties. For a random variable $X$ whose probability density function is:
  $$f(x) = \frac{1}{2} e^{-|x|}, \quad -\infty < x < \infty$$
  find the moment generating function of $X$ and hence find the mean and variance of $Z = 4X - 3$. \pts{7}
\end{parts}

\medskip\rule{\linewidth}{0.2pt}

\medskip


\noindent   \textbf{Question 5.}

\begin{parts}
  \item Let the joint probability mass function of random variables $(X, Y)$ be:
  $$P(X=x, Y=y) = \frac{x+2y}{18}, \quad x = 1, 2; \ y = 1, 2$$
  and zero elsewhere. Find:
  
\begin{parts}
    \item $E(X)$
    \item $E(Y|X)$
    \item $E[\max(X, Y)]$
  \end{parts} \pts{7}
  \item Define Bernoulli distribution. Let $X_1, X_2, \dots, X_n$ be $n$ independent Bernoulli random variables, each with parameter $p$. If $Y = \sum_{i=1}^n X_i$, obtain the probability distribution of $Y$. \pts{7}
\end{parts}

\medskip\rule{\linewidth}{0.2pt}

\medskip


\noindent   \textbf{Question 6.}

\begin{parts}
  \item Show that the mean, median, and mode for a normal distribution are always the same. \pts{7}
  \item The examination score ($X$) of 1000 students is normal with mean 50 and standard deviation 10. Approximately how many students will fail if the passing mark is 30? \pts{7}
\end{parts}

\medskip\rule{\linewidth}{0.2pt}

\medskip


\noindent   \textbf{Question 7.}

\begin{parts}
  \item Define Gamma distribution and find its mean, harmonic mean, and standard deviation. \pts{7}
  \item State and prove Chebyshev's inequality. If the mean of a random variable $X$ is zero and $P(|X| < 2) = 1/2$, find the lower bound for the variance of $X$. \pts{7}
\end{parts}

\medskip\rule{\linewidth}{0.2pt}

\medskip


\noindent   \textbf{Question 8.}

\begin{parts}
  \item Define the negative binomial distribution and find its mean. If the probability that a randomly selected person is found honest in interrogation is $\frac{1}{1000}$, how many persons are expected to be interrogated to get the first honest person? (Provide your logic for the answer). \pts{7}
  \item Write short notes on any four of the following in brief:
  
\begin{parts}
    \item Attributes and their analysis
    \item Coefficient of determination
    \item Bivariate normal distribution
    \item Intraclass correlation
    \item Poisson distribution
    \item Properties of normal distribution
  \end{parts} \pts{7}
\end{parts} 

\vfill
\rule{\linewidth}{0.4pt}

\begin{center}\small
  \href{https://bhu-exam.vercel.app/}{Click here to visit our website}\\[4pt]\href{https://me-aryan.vercel.app/}{Click here to visit our contributor}\\[4pt]\href{https://quashan-me.vercel.app/}{Click here to visit our contributor}
\end{center}

\end{document}